\section{Bounded formulas as signed finite games}

\subsection{Formula grammar}

Fix a finite formula schema with free parameter tuple \(\vec a\).  Its atoms
are reflected primitive-recursive relations.  The grammar is
\[
 A ::= R(\vec t)\mid A\wedge A\mid A\vee A
       \mid \exists x<b(\vec a,\vec z)\,A
       \mid \forall x<b(\vec a,\vec z)\,A.
\]
The bound may depend on parameters and variables chosen by earlier quantifiers,
but not on the variable it bounds.  The formula tree is fixed in the
metatheory.  All numerical bounds may vary with the input.

An atom is represented by a Boolean term \(r(\vec a)\in\{0,1\}\) and checked
derivations
\[
 r=1\to R,\qquad r=0\to\neg R.
\]
Equivalently one may supply \(r=1\leftrightarrow R\) together with Booleanity.
The negative direction is constructively available because the atom is
decidable.  Negation need not be pushed through the entire formula: the signed
compiler directly constructs a proof of \(\neg A\) from a negative strategy.

\subsection{Game reading}

Read a positive certificate as a strategy for the prover and a negative
certificate as a strategy for the refuter.

\begin{center}
\begin{tabular}{lll}
\toprule
Node & Positive certificate & Negative certificate \\
\midrule
Atom & recomputed true atom & recomputed false atom \\
\(A\wedge B\) & certificates for both & one failing conjunct \\
\(A\vee B\) & one successful disjunct & certificates for both \\
\(\exists x<b\,A_x\) & one witness and child certificate & table for every witness \\
\(\forall x<b\,A_x\) & table for every challenge & one counterexample and child certificate \\
\bottomrule
\end{tabular}
\end{center}

The table cases use heterogeneous mixed radix.  The one-child cases use
disjoint sums.  This duality is exactly the intuitionistic proof theory.  A
negative certificate for \(A\wedge B\) is not a proof of the classically
stronger formula \(\neg A\vee\neg B\) obtained by an unrestricted De Morgan
law; it explicitly contains the branch and a proof of that branch's negation.
Similarly, a negative certificate for a bounded existential explicitly
contains a refutation for every candidate.

\subsection{Value functions}

Define the Boolean value structurally:
\begin{align*}
 \Val_R &= r,\\
 \Val_{A\wedge B} &= \Val_A\cdot\Val_B,\\
 \Val_{A\vee B} &= \mathsf{or}(\Val_A,\Val_B),\\
 \Val_{\exists x<b A_x} &= \mathsf{any}_{x<b}\Val_{A_x},\\
 \Val_{\forall x<b A_x} &= \mathsf{all}_{x<b}\Val_{A_x}.
\end{align*}
Here \(\mathsf{all}\) is a finite product of normalized Booleans, and
\(\mathsf{any}\) may be defined as the sign of a finite sum or by a direct
primitive recursion.  Empty conventions are
\[
 \mathsf{any}_{x<0}=0,
 \qquad
 \mathsf{all}_{x<0}=1.
\]
Every value is primitive recursive, and Booleanity follows by structural
specialization and quantifier-free induction for the folds.

\subsection{Positive and negative code-space bounds}

Let \(P_A^+(\vec a)\) and \(P_A^-(\vec a)\) be the sizes of the syntactic
positive and negative code spaces.  The recurrences are
\begin{align*}
 P_R^+&=1,&P_R^-&=1,\\
 P_{A\wedge B}^+&=P_A^+P_B^+,&
 P_{A\wedge B}^-&=P_A^-+P_B^-,\\
 P_{A\vee B}^+&=P_A^++P_B^+,&
 P_{A\vee B}^-&=P_A^-P_B^-,\\
 P_{\exists x<b A_x}^+&=\sum_{x<b}P_{A_x}^+,&
 P_{\exists x<b A_x}^-&=\prod_{x<b}P_{A_x}^-,\\
 P_{\forall x<b A_x}^+&=\prod_{x<b}P_{A_x}^+,&
 P_{\forall x<b A_x}^-&=\sum_{x<b}P_{A_x}^-.
\end{align*}
A bound may be zero.  For example, the positive space of an empty existential
has size zero, while its negative space is the singleton empty table.  These
are exact bounds; they do not assert that every syntactically shaped code is
valid.

The sum/product pattern is the numerical form of the game duality.  One branch
is a disjoint sum; all branches are a Cartesian product encoded by mixed radix.
This observation provides a simple audit: if a constructor's code-bound
recurrence does not match its proof rule, the compiler is wrong.

\subsection{Checkers}

For each formula define primitive-recursive predicates
\[
 C_A^+(\vec a,c),\qquad C_A^-(\vec a,c)
\]
with values in \(\{0,1\}\).  Every checker first verifies
\(c<P_A^\pm(\vec a)\).

At an atom the only code is zero, and the checker recomputes the atomic Boolean.
At a product node, mixed-radix decoding produces the child codes and the
checker conjoins their checks.  At a disjoint node, prefix sums locate the owner
interval and the checker invokes the corresponding child.  At a bounded
quantifier, the same operations occur over a variable number of children;
the checker uses a finite Boolean fold.

The checker is intentionally independent of the canonical producer.  An
external algorithm may provide a different valid strategy.  Soundness trusts
only the checker equation, not the history of the code.

\subsection{Canonical codes}

The canonical producer follows the value.

\begin{itemize}
\item At a product node, recursively produce all child codes and pack them.
\item At a disjoint node, use bounded least search to find the least child with
the required value, recursively produce its code, and add the corresponding
prefix sum.
\item At an empty product, return zero.
\end{itemize}

There are two canonical terms \(\kappa_A^+\) and \(\kappa_A^-\), but only the
one matching \(\Val_A\) is required to be accepted.  Defining both as total
primitive-recursive functions avoids partial terms.  Their behavior on the
wrong side is irrelevant.

\subsection{Counterstrategy content}

A negative certificate is more informative than a proof of decidability.
For a formula
\[
 \forall q<Q\;\exists z<Z(q)\;\forall e<E(q,z)\;R(q,z,e),
\]
a positive certificate is a table assigning each \(q\) one \(z\) and, for
that \(z\), a table of atomic successes for all \(e\).  A negative certificate
chooses one \(q_0\) and contains, for every candidate \(z<Z(q_0)\), one
counterexample \(e_z\) with a negative atomic certificate.  This is exactly the
failure explanation practitioners want from a finite search.

No excluded-middle principle is hidden here.  The alternatives are computed by
finite primitive recursion over decidable atomic tests.  The object proof of
the winning side is generated together with the computation.
